\clearpage
\fancyhead[L]{\small JUNO.}
\hypertarget{native-255-lower}{}
\begin{center}
\textbf{ÜBER EINEN NEUEN VOM INSPECTOR HARDING IN LILIENTHAL ENTDECKTEN HÖCHST
MERKWÜRDIGEN WANDEL-STERN.}
\end{center}

\begin{center}
Monatliche Correspondenz. Band X. S. 371–385 1804 October.
\end{center}

\begin{center}
— — —
\end{center}

Kaum hatte ich dem unermüdlichen und unvergleichlichen Dr. \textsc{Gauss} meine drei ersten
Beobachtungen dieses Gestirns vom 13. 14. und 15. Sept. mitgetheilt, als ich mit
umgehender Post den 23. Sept. schon folgende Antwort erhielt, welche alle unsere Leser
eben so sehr, als mich, in Erstaunen setzen wird. „Was werden Sie sagen“, schreibt
dieser tiefsinnige Geometer, „dass ich es gewagt habe, auf meine eigenen Beobachtungen,
in Verbindung mit den drei mir von Ihrer Güte mitgetheilten und ein paar früheren
von Dr. \textsc{Olbers}, die zusammen nur eine Zeit von 14 Tagen und einen heliocentrischen
Bogen von vier Graden befassen, dass ich es gewagt habe, auf diese schlüpfrigen
Hülfsmittel schon einen vorläufigen Versuch und elliptische Elemente einer Bahn ohne
alle hypothetische Voraussetzungen zu gründen? Das Resultat kann nicht anders, als
sehr \textit{precär} sein; doch bin ich geneigt zu hoffen, dass es nicht mehr
\textit{enorm} oder \textit{total} von der Wahrheit abweichen kann, sondern wenigstens
schon einen rohen Begriff von den Dimensionen der Bahn gibt. Mit noch mehr Zuversicht
schmeichle ich mir, dass es zureichen wird, um allenfalls einen Monat hindurch,
vielleicht noch länger, den Planeten darnach aufzufinden; und mit Gewissheit kann ich
behaupten, dass alle bisherigen Beobachtungen gut dadurch dargestellt werden. Hier
einstweilen das Resultat, nächstens die Vergleichung mit den Beobachtungen, wobei ich
dann zugleich bestimmen werde, ob ich es des Titels: \textit{Elemente I des
\textsc{Harding}'schen Planeten}, würdig erklären kann.

\begin{center}
\begingroup\fontsize{9}{11}\selectfont\setlength{\tabcolsep}{3pt}
\begin{tabular}{lr}
Epoche Seeberger Merid. 1804 Sept. 5  & $24^\circ\ 53'\ 44''$ \\
Sonnenferne . . . . . . . . . . & $244\ 51\ 36$ \\
aufsteigender Knoten . . . . . . & $171\ 48\ 24$ \\
Excentricität . . . . . . . . . & $0.313757$ \\
halbe grosse Axe . . . . . . . . & $2.88208$ \\
tägliche Bewegung . . . . . . . & $725''18$ \\
Neigung der Bahn . . . . . . . & $15^\circ\ 12'\ 39''$ \\
Bewegung . . . . . . . . . . & rechtläufig \\
\end{tabular}\endgroup
\end{center}

„Was sagen Sie zu dieser sonderbaren Bahn, der grossen Excentricität, der grossen
Annäherung zur Ceres und Pallas, in Ansehung der Achse und mittlern Bewegung,
\textit{die gar leicht durch eine kleine Aenderung der Beobachtung zur völligen
Gleichheit werden kann}. Ich will aber meinem Grundsatze treu bleiben, den Rechnungen
schlechterdings nichts hypothetisches beizumischen, und künftigen Erfahrun%
\SourceLeaf{256}%
gen nicht vorzugreifen. In sehr kurzer Zeit werden wir schon viel weiter sein. Dass die
Bahn himmelweit von einer Parabel verschieden sei, und \textsc{Harding}'s Stern den
Planeten-Namen verdiene, daran lässt sich nun schon kaum mehr zweifeln; es wäre
daher zu wünschen, dass ihm bald ein Name beigelegt würde, natürlich muss das
Baptisations-Recht dem Entdecker allein vorbehalten bleiben u.\,s.\,w.

Hier sind die Beobachtungen, welche Dr. \textsc{Gauss} in Braunschweig angestellt hat.

\begin{center}
\begingroup\fontsize{9}{11}\selectfont\setlength{\tabcolsep}{3pt}\renewcommand{\arraystretch}{1.12}
\begin{tabular}{r|r|r|r}
1804 & \shortstack{Mittl. Zeit in\\Braunschweig} & \shortstack{Scheinbare ge-\\rade Aufsteig.} & \shortstack{Scheinbare\\südl. Declin.} \\ \hline
Sept. 12 & $10^{u}\,35^{m}\,2^{s}$ & $0^\circ\ 54'\ 26''$ & $1^\circ\ 38'\ 15''$ \\
13 & $9\ 41\ 32$ & $0\ 45\ 24$ & $1\ 50\ 59$ \\
14 & $11\ 38\ 49$ & $0\ 35\ 37$ & $2\ 5\ 1$ \\
15 & $10\ 16\ 17$ & $0\ 26\ 53$ & $2\ 17\ 42$ \\
16 & $10\ 37\ 4$ & $0\ 17\ 17$ & $2\ 31\ 20$ \\
17 & $11\ 28\ 59$ & $0\ 7\ 23$ & $2\ 44\ 7$ \\
\end{tabular}\endgroup
\end{center}

\begin{center}
\begingroup\fontsize{9}{11}\selectfont\setlength{\tabcolsep}{3pt}\renewcommand{\arraystretch}{1.12}
\begin{tabular}{r|r|r|r}
1804 & \shortstack{Mittl. Zeit in\\Braunschweig} & \shortstack{Scheinbare ge-\\rade Aufsteig.} & \shortstack{Scheinbare\\südl. Declin.} \\ \hline
18 & $11^{u}\,22^{m}\,16^{s}$ & $359^\circ\ 57'\ 26''$ & $2^\circ\ 58'\ 3''$ \\
21 & $10\ 24\ 52$ & $359\ 28\ 6$ & $3\ 37\ 59$ \\
24 & $10\ 1\ 2$ & $358\ 57\ 53$ & $4\ 18\ 6$ \\
25 & $8\ 44\ 25$ & $358\ 48\ 12$ & $4\ 30\ 44$ \\
27 & $10\ 20\ 29$ & $358\ 27\ 20$ & $4\ 57\ 47$ \\
28 & $8\ 29\ 4$ & $358\ 18\ 20$ & $5\ 10\ 22$ \\
\end{tabular}\endgroup
\end{center}

Er schreibt dazu: „Die drei ersten Beobachtungen sind mit einem schlechten, und
besonders schlecht montirten Achromaten gemacht, und verdienen daher wenig Vertrauen;
die nachfolgenden hingegen mit einem sehr guten Spiegel-Teleskop; diese werden daher
besser sein, wenigstens so gut, als es die Kreis-Mikrometer-Methode und meine
Gesichtsschärfe zulässt.“

Wenige Tage nach diesem Schreiben erhielt ich schon den 30. September die nähere
Bestätigung der Elemente dieser Planetenbahn. „Ich schicke Ihnen hier“, schreibt
Dr. \textsc{Gauss}, „neue und verbesserte Elemente, und schmeichle mir, dass sie schon eine
genäherte Bestimmung der wahren abgeben können, und wage es daher, sie als die
\textit{I. Elemente} anzukündigen.

\begin{center}
\begingroup\fontsize{9}{11}\selectfont\setlength{\tabcolsep}{3pt}
\begin{tabular}{lr}
Epoche Seeberg. Merid. 1804 5. Sept. . . . . . . . & $20^\circ\ 38'\ 56''$ \\
Sonnenferne . . . . . . . . . . . . . . . . & $239\ 14\ 2$ \\
aufsteigender Knoten . . . . . . . . . . . . . & $171\ 15\ 35$ \\
Excentricität . . . . . . . . . . . . . . . . & $0.287359$ \\
halbe grosse Axe (Logarithm) . . . . . . . . . . & $0.438682$ \\
tägliche Bewegung . . . . . . . . . . . . . . & $779''80$ \\
Neigung der Bahn . . . . . . . . . . . . . . & $13^\circ\ 34'\ 59''$ \\
Bewegung . . . . . . . . . . . . . . . . . & rechtläufig. \\
\end{tabular}\endgroup
\end{center}

Mit diesen neuen Elementen verglich nun Dr. \textsc{Gauss} die sämmtlichen Seeberger, Bremer
und Braunschweiger Beobachtungen; die Differenzen sind eben nicht grösser, als die
bei Kreis-Mikrometern möglichen Fehler, die zum Theil auch auf Rechnung der
verglichenen Sterne kommen mögen.

\begin{center}
\textit{Seeberger Beobachtungen:}
\end{center}

\begin{center}
\begingroup\fontsize{8}{10}\selectfont\setlength{\tabcolsep}{3pt}\renewcommand{\arraystretch}{1.12}
\begin{tabular}{r|r|r|r|r|r}
1804 & \shortstack{mittl. Zeit\\auf Seeberg} & \shortstack{Berechnete\\AR.} & \shortstack{Differenz\\R—B} & \shortstack{Berechnete\\Abweichung} & \shortstack{Differ.\\R—B} \\ \hline
Sept. 13 & $12^{u}\,31^{m}\,59^{s}$ & $0^\circ\ 44'\ 57''5$ & $+1''0$ & $1^\circ\ 52'\ 34''7$ & $-2''3$ \\
14 & $12\ 27\ 27$ & $0\ 35\ 43.8$ & $+2.8$ & $2\ 5\ 35.3$ & $-0.3$ \\
15 & $12\ 22\ 53$ & $0\ 26\ 20.4$ & $+0.4$ & $2\ 18\ 40.5$ & $+2.0$ \\
17 & $12\ 13\ 45$ & $0\ 7\ 7.8$ & $-0.3$ & $2\ 45\ 8.7$ & $-4.4$ \\
18 & $12\ 9\ 10$ & $359\ 57\ 20.4$ & $0$ & $2\ 58\ 28.3$ & $-2.6$ \\
20 & $ .\quad .\quad . $ & $359\ 37\ 28.7$ & $ .\ .\ . $ & $3\ 25\ 14.9$ & $+11.4$ \\
23 & $11\ 46\ 11$ & $359\ 7\ 11.9$ & $-3.0$ & $4\ 5\ 32.0$ & $-2.5$ \\
\end{tabular}\endgroup
\end{center}
\Needspace{23\baselineskip}
\SourceLeaf{257}%
\begin{center}
\textit{Bremer Beobachtungen:}
\end{center}

\begin{center}
\begingroup\fontsize{8}{10}\selectfont\setlength{\tabcolsep}{3pt}\renewcommand{\arraystretch}{1.12}
\begin{tabular}{r|r|r|r|r|r}
1804 & \shortstack{Mittl. Zeit\\in Bremen} & \shortstack{Berechnete\\AR.} & \shortstack{Differ.\\R—B} & \shortstack{Berechnete\\Abweichung} & \shortstack{Differ.\\R—B} \\ \hline
Sept. 7 & $10^{u}\,37^{m}\,21^{s}$ & $1^\circ\ 36'\ 57''$ & $+5''$ & $0^\circ\ 36'\ 6''$ & $-3''$ \\
8 & $8\ 11\ 20$ & $1\ 29\ 42$ & $+16$ & $0\ 47\ 4$ & $-15$ \\
9 & $10\ 48\ 50$ & $1\ 20\ 36$ & $+10$ & $1\ 0\ 47$ & $-3$ \\
10 & $8\ 15\ 6$ & $1\ 13\ 5$ & $+5$ & $1\ 11\ 58$ & $+2$ \\
11 & $10\ 34\ 3$ & $1\ 3\ 35$ & $+11$ & $1\ 25\ 51$ & $+10$ \\
12 & $11\ 18\ 32$ & $0\ 54\ 27$ & $+22$ & $1\ 39\ 3$ & $-1$ \\
13 & $8\ 54\ 0$ & $0\ 46\ 17$ & $+14$ & $1\ 50\ 41$ & $-9$ \\
14 & $8\ 24\ 44$ & $0\ 37\ 15$ & $+8$ & $2\ 3\ 27$ & $ .\quad . $ \\
15 & $10\ 54\ 28$ & $0\ 26\ 52$ & $+12$ & $2\ 17\ 57$ & $-8$ \\
17 & $10\ 23\ 9$ & $0\ 7\ 50$ & $+25$ & $2\ 44\ 12$ & $-20$ \\
18 & $8\ 38\ 17$ & $359\ 58\ 44$ & $-3$ & $2\ 56\ 35$ & $-16$ \\
21 & $8\ 30\ 54$ & $359\ 28\ 49$ & $-3$ & $3\ 36\ 50$ & $-4$ \\
21 & $10\ 9\ 32$ & $359\ 28\ 8$ & $-1$ & $3\ 37\ 45$ & $-1$ \\
23 & $13\ 25\ 57$ & $359\ 6\ 26$ & $+8$ & $4\ 6\ 32$ & $-5$ \\
24 & $8\ 27\ 37$ & $358\ 58\ 22$ & $+8$ & $4\ 17\ 12$ & $-10$ \\
25 & $8\ 41\ 38$ & $358\ 48\ 15$ & $+2$ & $4\ 30\ 32$ & $-22$ \\
\end{tabular}\endgroup
\end{center}

\begin{center}
\textit{Braunschweiger Beobachtungen;} 
\end{center}

\begin{center}
\begingroup\fontsize{8}{10}\selectfont\setlength{\tabcolsep}{3pt}\renewcommand{\arraystretch}{1.12}
\begin{tabular}{r|r|r|r|r|r}
1804 & \shortstack{mittl. Zeit in\\Braunschweig} & \shortstack{Berechnete\\AR.} & \shortstack{Differ.\\R—B} & \shortstack{Berechnete\\Abweichung} & \shortstack{Differ.\\R—B} \\ \hline
Sept. 12 & $10^{u}\,35^{m}\,2^{s}$ & $0^\circ\ 54'\ 46''$ & $+15''$ & $1^\circ\ 38'\ 36''$ & $+21''$ \\
13 & $9\ 41\ 32$ & $0\ 46\ 2$ & $+38$ & $1\ 51\ 2$ & $+3$ \\
14 & $11\ 38\ 49$ & $0\ 35\ 58$ & $+21$ & $2\ 5\ 15$ & $+14$ \\
15 & $10\ 16\ 17$ & $0\ 27\ 10$ & $+17$ & $2\ 17\ 33$ & $-2$ \\
16 & $10\ 37\ 4$ & $0\ 17\ 28$ & $+11$ & $2\ 30\ 58$ & $-22$ \\
17 & $11\ 28\ 59$ & $0\ 7\ 25$ & $+2$ & $2\ 44\ 44$ & $+15$ \\
18 & $11\ 22\ 16$ & $359\ 57\ 38$ & $+12$ & $2\ 58\ 4$ & $+1$ \\
21 & $10\ 24\ 52$ & $359\ 28\ 5$ & $-1$ & $3\ 37\ 50$ & $-10$ \\
24 & $10\ 1\ 2$ & $358\ 57\ 45$ & $-8$ & $4\ 18\ 1$ & $-5$ \\
25 & $8\ 44\ 25$ & $358\ 48\ 6$ & $-6$ & $4\ 30\ 43$ & $-1$ \\
27 & $10\ 29\ 29$ & $358\ 27\ 8$ & $-12$ & $4\ 58\ 17$ & $+30$ \\
28 & $8\ 29\ 4$ & $358\ 17\ 50$ & $-30$ & $5\ 10\ 30$ & $+8$ \\
\end{tabular}\endgroup
\end{center}

Auch Dr. \textsc{Olbers} hatte die Güte, uns seine fortgesetzten Beobachtungen des
merkwürdigen Fremdlings mitzutheilen:

\begin{center}
\begingroup\fontsize{9}{11}\selectfont\setlength{\tabcolsep}{3pt}\renewcommand{\arraystretch}{1.12}
\begin{tabular}{r|r|r|r}
1804 & \shortstack{Mittlere Zeit\\in\\Bremen} & \shortstack{Scheinbare\\gerade Auf-\\steigung} & \shortstack{Scheinbare\\südliche Ab-\\weichung} \\ \hline
Sept. 9 & $10^{u}\,48^{m}\,50^{s}$ & $1^\circ\ 20'\ 30''$ & $1^\circ\ 1'\ 5''$ \\
10 & $8\ 15\ 6$ & $1\ 12\ 55$ & $1\ 11\ 55$ \\
11 & $10\ 43\ 54$ & $1\ 3\ 20$ & $1\ 25\ 48$ \\
\end{tabular}\endgroup
\end{center}

Dr. \textsc{Olbers} schreibt dabei: „Diese Beobachtungen bedürfen noch einer kleinen
Verbesserung, da alle Vergleichungen, eine einzige ausgenommen, mit kleinen Sternen der
\textit{Hist. cél.} haben geschehen müssen. Da \textsc{Harding}'s und meine Beobachtung
wahrscheinlich die ersten sind, die man über diesen neuen Planeten angestellt hat, so
wäre es sehr zu wünschen, dass Sie die Gewogenheit hätten, die genauere
Bestimmung dieser kleinen \textsc{La Lande}'schen Sterne zu machen. Die Sterne sind folgende:

\Needspace{11\baselineskip}
\begin{center}\fontsize{9}{11}\selectfont
\begin{tabular}{l}
S. 119 Hist. cél.\\
$8$ Grösse\quad dritter Fad. $0^u\ 8^m\ 53^s5$\\
\quad (soll sein $0^u\ 7^m\ 53^s5$)\quad Zen. Dist. $49^\circ\ 20'\ 6''$
\end{tabular}
\[
\text{S. 131 Hist. cél.}\left\{\begin{array}{lll}
8\ \text{—}&\text{mittl. Fad.}\ 0\ 3\ 40.5&\text{— —}\ 49\ 59\ 0\\
8\ \text{—}&\text{— —}\ 0\ 4\ 26.0&\text{— —}\ 50\ 15\ 12
\end{array}\right.
\]
\end{center}
\par
\SourceLeaf{258}%
\indent Diese drei Sterne habe ich auch so genau, als möglich, bestimmt, und es folgen hier ihre mittleren Positionen für den Anfang des Jahres 1804:

\begin{center}
\begingroup\fontsize{8}{10}\selectfont\setlength{\tabcolsep}{3pt}\renewcommand{\arraystretch}{1.12}
\begin{tabular}{r|r|r|r|r}
Grösse & \shortstack{Mittlere gerade\\Aufsteigung\\1804} & \shortstack{Jährliche\\Verände-\\rung} & \shortstack{Mittlere südliche\\Abweichung\\1804} & \shortstack{Jährliche\\Verände-\\rung} \\ \hline
$8$ & $1^\circ\ 0'\ 24''0$ & $+46''0$ & $1^\circ\ 7'\ 19''6$ & $-20''0$ \\
$8$ & $1\ 11\ 52.7$ & $+46.0$ & $1\ 23\ 30.9$ & $-20.0$ \\
$8$ & $2\ 0\ 20.9$ & $+46.0$ & $0\ 28\ 15.7$ & $-20.0$ \\
\end{tabular}\endgroup
\end{center}

Es ist in der Geschichte der Astronomie aller Zeiten und aller Nationen beispiellos und es zeigt von der glänzenden Epoche der heutigen Sternkunde in Deutschland, dass ein Planet vorherverkündigt und in dem kurzen Zeitraum von drei Wochen zugleich entdeckt, beobachtet, seine Bahn berechnet und sein künftiger Lauf vorgezeichnet worden sei. Dies alles geschah jedoch durch die vereinten Kräfte vier Deutscher Astronomen, welche alles dieses schon geleistet hatten, ehe noch die Nachricht von der Existenz dieses neuen Weltkörpers unsere eifersüchtigen Nachbarn erreicht hatte.

Hier also zum Schluss eine Ephemeride des künftigen geocentr. Laufes dieses neuen Planeten, welche Dr. \textsc{Gauss} nach seinen obigen I. Elementen berechnet hat. So grosser Verbesserungen auch diese Elemente noch bedürfen mögen, so hofft Dr. \textsc{Gauss} doch mit Zuversicht, dass sie mehrere Wochen hindurch diesen Lauf genau darstellen werden: daher wird folgende Ephemeride denjenigen entfernten Astronomen sehr willkommen sein, welche diesen neuen Himmels-Gast noch nicht aufgefunden haben. Die Momente sind für Mitternacht in Seeberg gerechnet.

\begin{center}
\begingroup\fontsize{9}{11}\selectfont\setlength{\tabcolsep}{3pt}\renewcommand{\arraystretch}{1.12}
\begin{tabular}{r|r|r}
\shortstack{Mitter-\\nacht in\\Seeberg} & \shortstack{Gerade Aufst.\\des neuen\\Planeten} & \shortstack{Südl. Abw.\\des neuen\\Planeten.} \\ \hline
Sept. 5 & $1^\circ\ 52'$ & $0^\circ\ 12'$ \\
12 & $0\ 54$ & $1\ 39$ \\
19 & $359\ 47$ & $3\ 12$ \\
26 & $358\ 36$ & $4\ 55$ \\
30 & $357\ 56$ & $5\ 39$ \\
Octob. 3 & $357\ 27$ & $6\ 17$ \\
6 & $357\ 0$ & $6\ 54$ \\
9 & $356\ 35$ & $7\ 29$ \\
12 & $356\ 12$ & $8\ 2$ \\
15 & $355\ 51$ & $8\ 33$ \\
\end{tabular}\endgroup
\end{center}
\begin{center}
\begingroup\fontsize{9}{11}\selectfont\setlength{\tabcolsep}{3pt}\renewcommand{\arraystretch}{1.12}
\begin{tabular}{r|r|r}
\shortstack{Mitter-\\nacht in\\Seeberg} & \shortstack{Gerade Aufst.\\des neuen\\Planeten} & \shortstack{Südl. Abw.\\des neuen\\Planeten} \\ \hline
Octob. 18 & $355^\circ\ 35'$ & $9^\circ\ 1'$ \\
21 & $355\ 21$ & $9\ 27$ \\
24 & $355\ 11$ & $9\ 50$ \\
27 & $355\ 4$ & $10\ 11$ \\
30 & $355\ 2$ & $10\ 28$ \\
Nov. 2 & $355\ 4$ & $10\ 42$ \\
5 & $355\ 10$ & $10\ 54$ \\
8 & $355\ 20$ & $11\ 3$ \\
11 & $355\ 34$ & $11\ 8$ \\
14 & $355\ 53$ & $11\ 13$ \\
\end{tabular}\endgroup
\end{center}

\hfill[\textsc{von Zach.}]
